\documentclass[11pt]{article}
\usepackage[margin=1.15in]{geometry}
\usepackage{amsmath,amssymb,amsthm}
\usepackage{booktabs}
\usepackage[hidelinks]{hyperref}

\newtheorem{remark}{Remark}
\newcommand{\eps}{\varepsilon}
\newcommand{\loglog}{\log\log}

\title{Rubinstein--Sarnak densities for weighted prime number races}
\author{Raghav Sharma\thanks{Draft of August 2026, not yet submitted.
This work was carried out with substantial AI assistance (Anthropic's
Claude): derivations, implementations, and the literature review were
AI-produced and then machine- and hand-verified as described in
Section~\ref{sec:methods}; the author takes responsibility for the
content. Comments are very welcome. Contact: raghavsharma1729@gmail.com.
Preprint: \href{https://doi.org/10.5281/zenodo.21947102}{doi:10.5281/zenodo.21947102};
code and data: \url{https://github.com/raghavs1729/weighted-prime-races}.}}
\date{August 2026}

\begin{document}
\maketitle

\begin{abstract}
For $m\in\mathbb R$ consider the weighted prime race
$D_m(x)=\sum_{p\le x,\,p\equiv 3\,(4)}p^{\,m}-\sum_{p\le x,\,p\equiv
1\,(4)}p^{\,m}$ and its analogues for other moduli. The regime
$m\le 0$ has a developed literature: the critical
$\tfrac12\loglog x$ deflection at $m=-\tfrac12$ (Aoki--Koyama; Sheth),
the sign-change threshold (Shimada--Koyama), and, for $m<-\tfrac12$,
limits expressible through Meissel--Mertens constants for progressions
(Languasco--Zaccagnini). We contribute the pieces this literature does
not contain: (i) the first numerical Rubinstein--Sarnak logarithmic
densities for a weighted race between primes in arithmetic progressions
(for weighted prime-sum races against main terms, cf.\ Lamzouri's
Mertens-race density~\cite{Lam16}), computed across the delicate
range under the standard GRH + LI hypotheses by deterministic Fourier
inversion --- e.g.\ for modulus $4$, $\delta(1)=0.79763$,
$\delta(2)=0.69452$, $\delta(3)=0.64585$, with the calibration
$\delta(0)=0.995928$ recovering the classical Rubinstein--Sarnak value
to six decimals; (ii) the
growing-weight regime $m>0$, previously untouched, with exact sieve
computations to $10^{10}$ agreeing with the predicted densities to
$\lesssim4\times10^{-3}$ for modulus $4$; (iii) an asymptotic dissolution law in the weight
aspect, $\delta_q(m)-\tfrac12\sim(2\pi\cdot
2(m+\tfrac12)\log\tfrac{q(m+1/2)}{2\pi})^{-1/2}$, an analogue of the
Fiorilli--Martin density asymptotics, verified numerically in moduli
$3$ and $4$. Taken together these results connect the weighted
formulation of Chebyshev's bias with the classical density formulation,
contributing computational evidence toward Problem 12 of the comparative
prime number theory problem list of Hamieh, Kadiri, Martin and Ng (the
logical relation the problem asks for remains a theoretical question).
\end{abstract}

\section{Introduction}

Chebyshev's observation --- that primes $\equiv3\ (\mathrm{mod}\ 4)$
predominate over primes $\equiv1\ (\mathrm{mod}\ 4)$ for most $x$ ---
was quantified by Rubinstein and Sarnak~\cite{RS94}: under GRH and the
linear independence (LI) of the ordinates of zeros of $L(s,\chi_4)$,
the set $\{x:\pi(x;4,3)>\pi(x;4,1)\}$ has logarithmic density
$0.9959\ldots$. A separate line of work, initiated by Aoki and
Koyama~\cite{AK23} and developed by Sheth~\cite{Sheth24} and
Shimada--Koyama~\cite{SK25}, studies \emph{weighted} counts
$\sum_{p\le x,\,p\equiv a\,(q)}p^{-s}$ and locates a transition at
$s=\tfrac12$: for the mod-$4$ race the weighted difference obeys
$\sum_{p\le x}\chi_4(p)p^{-1/2}=-\tfrac12\loglog x+O(1)$ (under DRH in
\cite{AK23}, under GRH off a small exceptional set in \cite{Sheth24}),
so the weighted race is single-signed for large $x$, while for weights
$p^{-s}$ with $s<\tfrac12$ sign changes persist~\cite{SK25}; Hayani
\cite{Hay25} has recently obtained the density-one statement at
$s=\tfrac12$ under GRH alone, without linear independence. Humphries
\cite{Hum13} established the analogous phase structure for weighted
sums of the Liouville function, including the limit
$\delta_\alpha\to1$ as $\alpha\to\tfrac12^-$. Problem 12 of the
comparative prime number theory problem list~\cite{HKMN24} (posed by
Koyama) asks how the weighted formulation relates to the classical
logarithmic-density formulation.

This note treats the two formulations as endpoints of one family. For
$m\in\mathbb R$ set
\[
D_m(x)\;=\;\sum_{\substack{p\le x\\ p\equiv 3\,(4)}}p^{\,m}
-\sum_{\substack{p\le x\\ p\equiv 1\,(4)}}p^{\,m},
\]
so $m=0$ is the classical race and $m=-s$ the Aoki--Koyama weighting;
we allow $m>0$, which appears not to have been considered. Our results
are computational, at the standard of numerical work in this area
(GRH + LI where densities are involved), and every headline number was
computed by two independent routes.

\subsection*{Results}

\textbf{1. Densities across the delicate range.} For $m>-\tfrac12$ the
limiting logarithmic distribution of the normalized race exists by the
framework of Akbary--Ng--Shahabi~\cite{ANS14} (cf.\ Devin~\cite{Dev20})
applied to the shifted explicit formula; the normalized race has the
almost-periodic form $1+\sum_{\gamma>0}2a_\gamma(m)\cos(\gamma u)$ with
\[
a_\gamma(m)\;=\;\frac{2m+1}{\sqrt{(m+\tfrac12)^2+\gamma^2}},
\]
$\gamma$ the ordinates of $L(s,\chi_4)$, and the constant $1$
normalizing the prime-square bias term. Monte Carlo evaluation under
LI (Section~\ref{sec:methods}) gives Table~\ref{tab:delta}. These are,
to our knowledge, the first Rubinstein--Sarnak densities computed for a
weighted race between primes in arithmetic progressions (for weighted
prime sums against main terms, densities exist in the Mertens-race
literature~\cite{Lam16}; for analytic explicit-formula control of the
log-weighted sum $\theta(x)$, cf.\ Platt--Trudgian~\cite{PT14}).

\begin{table}[h]\centering
\begin{tabular}{lccccccc}
\toprule
$m$ & $0$ & $1$ & $2$ & $3$ & $8$ & $20$ & $100$\\
\midrule
$\delta_4(m)$ & $0.995928$ & $0.797628$ & $0.694520$ & $0.645850$
& $0.571869$ & $0.538455$ & $0.513778$\\
$\delta_3(m)$ & $0.999063$ & $0.836877$ & $0.723470$ & $0.666552$
& $0.578534$ & $0.540767$ & --- \\
\bottomrule
\end{tabular}
\caption{Logarithmic densities of the set where the nonresidue class
leads, moduli $4$ and $3$: deterministic Gil--Pelaez Fourier inversion
of the limiting distribution ($511$ resp.\ $537$ zeros computed to $25$
digits, Gaussian tail from the zero-counting density), cross-checked by
independent $2\times10^7$-sample Monte Carlo to $\le10^{-4}$. The $m=0$
column reproduces the classical values $0.995928$ and $0.999063$
(cf.\ \cite{RS94,FM13}) to all six decimals.}
\label{tab:delta}
\end{table}

\textbf{2. Exact computations, $m>0$.} By segmented sieve with exact
$128$-bit signed accumulators to $X=10^{10}$: the mod-$4$ race first
favours the residue class at $p=5$ for every $m\ge1$ (against Leech's
$26{,}861$ for $m=0$); lead changes to $10^{10}$ number $3{,}082$
($m=0$), $125{,}624$, $159{,}942$, $199{,}540$ ($m=1,2,3$); the
finite-$x$ logarithmic measures at $10^{10}$ are $0.79380$, $0.69275$,
$0.64900$ for $m=1,2,3$, within $\lesssim4\times10^{-3}$ of
Table~\ref{tab:delta}. For modulus $3$ the weighted first flips occur
at $p=13$ ($m=1$) and $p=7$ ($m=2,3$), while the classical race, as
expected, does not flip below $10^{10}$.

\textbf{3. Dissolution law.} Substituting the zero-counting density
$\frac{1}{2\pi}\log\frac{qt}{2\pi}$ into the variance of the limiting
distribution gives
\[
\mathrm{Var}_q(m)\;\sim\;2\bigl(m+\tfrac12\bigr)
\log\frac{q(m+\tfrac12)}{2\pi},\qquad
\delta_q(m)-\tfrac12\;\sim\;\frac{1}{\sqrt{2\pi\,\mathrm{Var}_q(m)}}
\quad(m\to\infty),
\]
a weight-aspect analogue of the Fiorilli--Martin density
asymptotics~\cite{FM13} (cf.\ Meng~\cite{Meng18} for the $k$-aspect).
At $q=4$, $m=100$: deterministic $0.513778$, Gaussian $\Phi(1/\sigma)$
$0.51379$, the law $0.51380$. At $q=3$, $m=60$: deterministic $0.519716$
against the law's $0.51978$ --- agreement to $6\times10^{-5}$. The modulus ordering
$\delta_3(m)>\delta_4(m)$ predicted by the law (through
$\gamma_1(\chi_3)=8.04>\gamma_1(\chi_4)=6.02$) holds at every computed
$m$.

\textbf{4. The boundary $m\le-\tfrac12$, numerically.} Confirming (not
discovering) the theorems cited above: with
$C^*=\tfrac12(M-\log 2-\tfrac12)+K_3-\log L(\tfrac12,\chi_4)
=-0.07153$ ($M$ the Meissel--Mertens constant, $K_3$ the $k\ge3$
prime-power contribution), the parameter-free prediction
$D_{-1/2}(x)=\tfrac12\loglog x+C^*+\eps(x)$ gives $1.4968\pm0.0172$
at $x=10^{10}$ against the exact value $1.50775$; the fitted slope
over $[10^6,10^{10}]$ is $0.494$. The fluctuation $\eps$ decays like
$\sqrt{\Sigma}/\log x$ with $\Sigma=\sum_\gamma2/\gamma^2=0.156033$,
visible decade-by-decade in the residuals --- a small refinement of
the qualitative picture. For $m<-\tfrac12$ the limits are identified
with differences of Meissel--Mertens constants for progressions:
$\lim D_{-1}=M(4,3)-M(4,1)=0.334981325299993\ldots$, our
M\"obius-over-$\log L$ evaluation agreeing with
Languasco--Zaccagnini~\cite{LZ10} to all compared digits (the
difference itself is OEIS A086239; the individual constants are
A368645 and A368646). Notably, none of the races $m\in\{-1,-\tfrac34,
-\tfrac12\}$ changes sign even once on $[3,10^{10}]$ (minimum margin
$1/\sqrt3-1/\sqrt5\approx0.13014$, attained by the $m=-\tfrac12$ race
at $x=5$; the $m=-1$ margin there is $2/15$): in this range of weights Chebyshev's bias is not
merely predominant but, empirically, total from the first odd prime.

\section{Methods and verification}\label{sec:methods}

Zeros of $L(s,\chi_4)$ and $L(s,\chi_3)$ ($511$ resp.\ $537$ ordinates
to $25$ digits) were computed by Hardy $Z$-scans with derived phases,
validated by reality of $Z$ to $\sim10^{-27}$, Riemann--von Mangoldt
counts, and the known first ordinates. Densities: Monte Carlo over
independent uniform phases (LI), $2$--$4\times10^6$ samples, with the
unseen-zero tail added as an independent Gaussian with variance from
the counting density; calibration anchors $\delta_4(0)=0.99593$,
$\delta_3(0)=0.99904$ against \cite{RS94}. Exact races: segmented
sieves to $10^{10}$ with exact signed $128$-bit (integer weights) or
compensated long-double (fractional weights) accumulation,
cross-checked against independent implementations at overlapping
ranges; the $m=0$ anchors (first flip $26{,}861$; no mod-$3$ flip
below $10^{10}$) reproduce the classical values. The density kernel is
derived for the $\log$-weighted ($\theta$-form) race; partial summation
gives every term of the unlogged race the same $1/\log x$ factor, which
cancels under the bias normalization, so the limiting distribution and
density coincide --- the same passage Rubinstein--Sarnak make at $m=0$. The explicit-formula
kernel derives from a uniform smoothed-Perron framework developed in a
companion manuscript (analytic computation of $S_m(x)$ from zeros,
exact to $10^9$ at $m=1$), available with all code and data at
\url{https://github.com/raghavs1729/weighted-prime-races} (archived at
\href{https://doi.org/10.5281/zenodo.21947102}{doi:10.5281/zenodo.21947102}). Every headline number in this note was
re-computed independently of the code that first produced it.

\emph{Assumptions and limitations.} Densities are conditional on
GRH + LI, as throughout this literature; the existence of the limiting
distribution for $m>-\tfrac12$ is a routine adaptation of
\cite{ANS14,Dev20} and is not claimed as new; the dissolution law is
derived at the level of rigor of \cite{FM13}'s heuristic term plus
numerical confirmation, and a fully rigorous proof (with error term)
is left open. The finite-$x$ log-measure for $(q,m)=(3,2)$ shows the
weakest agreement in the study ($0.023$ at $10^{10}$), plausibly a
slow almost-periodic phase; we flag it rather than explain it.

\begin{thebibliography}{99}
\bibitem{RS94} M. Rubinstein, P. Sarnak, \emph{Chebyshev's bias},
Experiment.\ Math.\ \textbf{3} (1994), 173--197.
\bibitem{AK23} M. Aoki, S. Koyama, \emph{Chebyshev's bias against
splitting and principal primes in global fields}, J. Number Theory
\textbf{245} (2023), 233--262. arXiv:2203.12266.
\bibitem{Sheth24} A. Sheth, \emph{Euler products at the centre and
applications to Chebyshev's bias}, Math.\ Proc.\ Cambridge Philos.\
Soc.\ \textbf{179} (2025), 331--349. arXiv:2405.01512.
\bibitem{SK25} K. Shimada, S. Koyama, \emph{Weighted prime number
theorem on arithmetic progressions with refinements}, Mathematics
\textbf{13} (2025), no.\ 16, 2564.
\bibitem{Hum13} P. Humphries, \emph{The distribution of weighted sums
of the Liouville function and P\'olya's conjecture}, J. Number Theory
\textbf{133} (2013), 545--582. arXiv:1108.1524.
\bibitem{HKMN24} A. Hamieh, H. Kadiri, G. Martin, N. Ng,
\emph{Comparative prime number theory problem list}, arXiv:2407.03530
(2024), Problem 12.
\bibitem{ANS14} A. Akbary, N. Ng, M. Shahabi, \emph{Limiting
distributions of the classical error terms of prime number theory},
Q. J. Math.\ \textbf{65} (2014), 743--780.
\bibitem{Dev20} L. Devin, \emph{Chebyshev's bias for analytic
$L$-functions}, Math.\ Proc.\ Cambridge Philos.\ Soc.\ \textbf{169}
(2020), 103--140. arXiv:1706.06394.
\bibitem{FM13} D. Fiorilli, G. Martin, \emph{Inequities in the
Shanks--R\'enyi prime number race: an asymptotic formula for the
densities}, J. reine angew.\ Math.\ \textbf{676} (2013), 121--212.
\bibitem{Meng18} X. Meng, \emph{Chebyshev's bias for products of $k$
primes}, Algebra Number Theory \textbf{12} (2018), 305--341.
\bibitem{LZ10} A. Languasco, A. Zaccagnini, \emph{Computing the
Mertens and Meissel--Mertens constants for sums over arithmetic
progressions}, Experiment.\ Math.\ \textbf{19} (2010), 279--284.
\bibitem{Fio14} D. Fiorilli, \emph{Highly biased prime number races},
Algebra Number Theory \textbf{8} (2014), 1733--1767.
\bibitem{Lam16} Y. Lamzouri, \emph{A bias in Mertens' product formula},
Int.\ J. Number Theory \textbf{12} (2016), 97--109. arXiv:1410.3777.
\bibitem{Hay25} M. Hayani, \emph{Chebyshev's bias without linear
independence}, arXiv:2512.23302 (2025).
\bibitem{PT14} D.~J. Platt, T.~S. Trudgian, \emph{On the first sign
change of $\theta(x)-x$}, Math.\ Comp.\ \textbf{85} (2016), 1539--1547.
arXiv:1407.1914.
\end{thebibliography}

\end{document}
